% !TEX program = xelatex
\documentclass{article}
\usepackage{geometry}\geometry{margin=23mm}
\usepackage{fontspec}\setmainfont{Times New Roman}
\usepackage[italian]{babel}
\begin{document}
\title{La tavola di misura}\author{}\date{14 maggio 2026}\maketitle
\begin{abstract}Due segni indicano la posizione iniziale della tavola.\end{abstract}
\tableofcontents
\section{I segni}\begin{quote}«Leggere prima il segno inferiore».\end{quote}
\begin{figure}[ht]\centering\fbox{\rule{0pt}{8mm}\rule{33mm}{0pt}}\caption{Posizione della tavola}\end{figure}
\section{Altezza}
\begin{table}[ht]\centering\begin{tabular}{cr}Segno&Altezza (mm)\\\hline 1&12\\2&18\end{tabular}\caption{Altezza dei segni}\end{table}
\begin{thebibliography}{1}\bibitem{record}\emph{Quaderno di misura}, 2025, \chaptername~2.\end{thebibliography}
\end{document}
